$27y^2=4(x-2)^3$ desde $(2,0)$ a $(11,6\sqrt{3})$.

\nopagebreak
\begin{align*}
y^2 = \frac{4}{27}(x-2)^3 \implies y = \frac{2}{3\sqrt{3}}(x-2)^{3/2}. \\
y' = \frac{2}{3\sqrt{3}} \cdot \frac{3}{2} (x-2)^{1/2} = \frac{1}{\sqrt{3}} \sqrt{x-2}. \\
1 + (y')^2 = 1 + \frac{1}{3}(x-2) = \frac{3+x-2}{3} = \frac{x+1}{3}. \\
L &= \int_{2}^{11} \sqrt{\frac{x+1}{3}} \, dx = \frac{1}{\sqrt{3}} \int_{2}^{11} (x+1)^{1/2} \, dx \\
&= \frac{1}{\sqrt{3}} \left[ \frac{2}{3} (x+1)^{3/2} \right]_{2}^{11} \\
&= \frac{2}{3\sqrt{3}} (12^{3/2} - 3^{3/2}) \\
&= \frac{2}{3\sqrt{3}} (12\sqrt{12} - 3\sqrt{3}) = \frac{2}{3\sqrt{3}} (24\sqrt{3} - 3\sqrt{3}) \\
&= \frac{2}{3\sqrt{3}} (21\sqrt{3}) = \frac{42}{3} = 14
\end{align*}

\[
\boxed{\displaystyle 14}
\]
